% ---
% Agradecimentos
% ---

\begin{agradecimentos}

Agradecemos ao Prof. Me. Ícaro de Paula Freitas, que orientou este trabalho em sua primeira metade, pela confiança no projeto e pelas orientações que deram a ele a sua forma inicial.

Agradecemos ao Prof. Me. Raphael Lopes de Souza, nosso orientador, pela disponibilidade, pela paciência e pelas orientações que conduziram este trabalho até a sua conclusão.

Agradecemos à Profa. Ma. Regina Albanese Pose, que participou das reuniões iniciais e, com o seu conhecimento em estatística e em tecnologia aplicada à saúde, nos guiou na direção de um projeto bem-sucedido.

Agradecemos de forma especial ao Prof. Dr. Álex Aparecido Rosini Silva e ao técnico de laboratório Marcos Eduardo Souza Abreu, do Laboratório de Metabolômica e Proteômica da USCS, vinculado ao INOVA-USCS, que se dispuseram a ser parte integrante deste projeto: apresentaram o problema, ajudaram a delimitar o escopo, esclareceram as nossas dúvidas sobre o assunto, recomendaram bibliografia, avaliaram a monografia e avaliaram o projeto final.

Agradecemos ao INOVA-USCS, centro de inovação da Universidade Municipal de São Caetano do Sul, por meio dos professores Dr. Carlos Eduardo Bognar e Dra. Sheila de Oliveira Garcia Mateos, que tornaram possível o contato com o Laboratório de Metabolômica e o acesso às suas necessidades.

Ao Prof. Me. Fabio Rezende de Souza, que demonstrou seu interesse no projeto e se prontificou a oferecer apoio caso precisássemos, nosso agradecimento pela generosidade e pela atenção.

Agradecemos também ao corpo docente da Universidade Municipal de São Caetano do Sul que cruzou o caminho das nossas turmas. O tempo e o esforço que esses professores dedicaram à nossa formação são o que nos permitiu chegar até aqui e realizar um projeto tão especial para nós.

\end{agradecimentos}

% ---

